\documentclass[11pt,a4paper,sans]{moderncv}

\moderncvstyle{classic}
\moderncvcolor{blue}

% first and last name with the same weight for a more balanced title
\renewcommand*{\namefont}{\fontsize{30}{32}\mdseries\upshape}
\definecolor{softaccent}{RGB}{77, 136, 204}
\newcommand{\company}[1]{\textcolor{softaccent}{\textbf{#1}}}
\newcommand{\metric}[1]{\textcolor{softaccent}{\textbf{#1}}}
\let\oldsection\section
\renewcommand{\section}[1]{\oldsection{#1}\vspace{0.25em}}
\let\oldcventry\cventry
\renewcommand*{\cventry}[6]{\oldcventry{#1}{#2}{#3}{#4}{#5}{#6}\vspace{0.35em}}
\newcommand{\customcvtitle}{%
  \begin{minipage}[c]{0.74\textwidth}
    {\namefont Muhammet Emin Balmuk\par}
    \vspace{0.25em}
    {\large\itshape Computer Engineering Student\par}
    \vspace{0.6em}
    {\small Maltepe/Istanbul, Turkey \textbar\ +90 533 365 6944 \textbar\ \href{mailto:eminblmk@gmail.com}{eminblmk@gmail.com}\par}
    \vspace{0.2em}
    {\small Website: \href{https://eminbalmuk.github.io}{eminbalmuk.github.io} \textbar\ LinkedIn: \href{https://www.linkedin.com/in/eminbalmuk}{eminbalmuk} \textbar\ GitHub: \href{https://github.com/eminbalmuk}{eminbalmuk}\par}
  \end{minipage}\hfill
  \begin{minipage}[c]{0.2\textwidth}
    \raggedleft
    \setlength{\fboxsep}{1.5pt}%
    \fcolorbox{gray!55}{white}{\includegraphics[width=0.93\linewidth]{photo.jpg}}
  \end{minipage}
  \par\vspace{0.8em}
  \rule{\textwidth}{0.4pt}
  \par\vspace{0.6em}
}
\newcommand{\stackitem}[2]{%
  \cvitem{#1}{#2}%
}
\newcommand{\skilldivider}{%
  \vspace{0.15\baselineskip}%
  \cvitem{}{\textcolor{gray!45}{\rule{\linewidth}{0.3pt}}}%
  \vspace{0.05\baselineskip}%
}
\newcommand{\entrygap}{\vspace{\baselineskip}}

\usepackage[scale=0.82, top=1.7cm, bottom=1.7cm]{geometry}
\usepackage[T1]{fontenc}
\usepackage[utf8]{inputenc}
\usepackage[english]{babel}
\usepackage{hyperref}
\usepackage{needspace}
\hypersetup{colorlinks=true, urlcolor=blue}

% personal information
\name{Muhammet Emin Balmuk}
\title{}
\address{Maltepe/Istanbul}{Turkey}
\phone[mobile]{+90 533 365 6944}
\email{eminblmk@gmail.com}
\social[linkedin]{eminbalmuk}
\social[github]{eminbalmuk}

% picture
\photo[64pt][0.4pt]{photo.jpg}

\begin{document}
\customcvtitle
\vspace{-0.4em}

\section{Profile}
I am a third-year Computer Engineering student at Marmara University. Through the 8-month data science training I completed at the Google AI and Technology Academy, I strengthened my skills in Python, data analysis, and machine learning. I aim to build my career in Software Development, Data Science, and AI Engineering by contributing to software, AI, machine learning, and data-driven projects with a collaborative mindset and a strong focus on continuous improvement.

\entrygap
\section{Education}
\cventry{2023--2027}{B.Sc. in Computer Engineering}{Marmara University}{Istanbul}{}{Third-year student}
\cventry{2019--2023}{High School}{Bah\c{c}elievler Anatolian High School (German)}{}{}{}

\section{Experience}
\cventry{December 2024--August 2025}{Data Science Trainee}{\company{Google AI and Technology Academy}}{Istanbul}{}{
\begin{itemize}
    \item Completed an intensive 8-month data science training program.
    \item Built hands-on knowledge in data analysis, statistics, machine learning, and project development.
  \item Emphasized model validation discipline and business-facing interpretation of technical outputs.
    \item Strengthened teamwork, model development, and problem-solving skills through datathon and bootcamp projects.
\end{itemize}}

\entrygap

\cventry{February 2026--Present}{ILERIDE Program \textbar\ Trainee}{\company{Shell}}{Istanbul}{}{
\begin{itemize}
    \item Participating in an ongoing development-focused program centered on structured thinking, problem framing, and professional business perspectives.
  \item Developing clearer decision-support communication by translating analytical work into business-understandable framing.
    \item Applying takeaways from program sessions and activities to better understand how corporate teams approach decision-making and problem solving.
\end{itemize}}

\entrygap

\cventry{January 2026--Present}{Project Development Team \textbar\ Core Member}{\company{GDG - Google Developer Group Marmara}}{}{}{
\begin{itemize}
    \item Support technical project initiatives as part of the GDG Marmara core team.
  \item Contribute to scoped delivery planning and sustainable technical execution in community-driven projects.
    \item Collaborate on project planning, product thinking, and software-focused community work.
\end{itemize}}

\newpage
\section{Projects}

\cventry{2025--Present}{Crypto Analysis and Trading Bot}{Personal Project}{}{}{
\begin{itemize}
    \item Engineered a Python-based algorithmic trading bot built on custom quantitative strategies.
    \item Built an architecture that screens \metric{200+ crypto assets} in about \metric{2 minutes} and reached a \metric{66\% success rate} with a \metric{1:2.5 risk-reward ratio} in a real-time Telegram signal system.
  \item Implemented a custom labeling workflow via Google Sheets API and integrated the outputs into the modeling process.
    \item Evaluated LightGBM, CatBoost, and TabNet during R\&D, then refactored the live system into a custom strategy-based architecture to reduce latency.
    \item Processed high-volume market data with NumPy and Pandas and maintained trade logs and historical records in MongoDB.
    \item Automated 24/7 trade lifecycle tracking through an additional ``Crypto Signal Tracker'' bot.
\end{itemize}}

\entrygap

\cventry{2025}{Datathon Competition}{Google AI and Technology Academy}{}{}{
\begin{itemize}
    \item Analyzed and visualized a \metric{210,000-row dataset} as part of a \metric{4-person team}.
    \item Built price prediction models for a new \metric{56,000-row dataset} using LightGBM, Random Forest, and XGBoost.
    \item Interpreted model outputs and turned findings into visual insights with Matplotlib and Seaborn.
\end{itemize}}

\entrygap

\cventry{2025}{AI-Integrated Health Application}{Bootcamp Project}{}{}{
\begin{itemize}
    \item Built a mobile application with a \metric{5-person team} for the preliminary assessment of cognitive conditions such as hyperactivity, dementia, and autism.
    \item Delivered Python backend and Dart frontend components; ranked in the \metric{top 3 among 205 teams}.
    \item Developed a Python voice pronunciation test module, managed Firebase operations, and integrated backend and frontend through REST APIs.
    \item Improved selected Dart screens and helped address application security issues.
    \item \href{https://github.com/BBBeyza/YZTA_YapayZeka_Grup-132}{GitHub Repository}
\end{itemize}}

  \entrygap

\cventry{2025}{Fire Detection System for Live Security Cameras}{Computer Vision}{}{}{
\begin{itemize}
    \item Trained a model on a dataset of \metric{95,000 images} using YOLO 11.
    \item Reached \metric{82.8\%} smoke detection accuracy and \metric{79.6\%} fire detection accuracy, with especially strong results on black-and-white images.
    \item Tested the system on both live camera streams and recorded video footage, marking detected regions with bounding boxes.
\end{itemize}}

  \entrygap

\cventry{2025}{Real-Time Chat Application}{Team Project}{}{}{
\begin{itemize}
    \item Built a real-time messaging application with a 3-person team using HTML, CSS (Tailwind), JavaScript, React, Node.js, MongoDB, and Socket.io.
    \item Managed Node.js backend tasks, real-time connection management with Socket.io, MongoDB workflows, and media storage with Cloudinary.
\end{itemize}}

  \entrygap

\cventry{2024}{Energy Distribution and Fault Response Simulation}{Team Project}{}{}{
\begin{itemize}
    \item Built a simulation project with a 4-person team using Java (Spring Boot), JavaScript, and HTML to model fault detection and response during a possible power outage in Istanbul.
    \item Built the full JavaScript/HTML interface and supported backend integration for the outage-response workflow.
\end{itemize}}

  \entrygap

\cventry{2024}{Personal Series/Movie Watchlist Platform}{Team Project}{}{}{
\begin{itemize}
    \item Developed an application with a 3-person team using Python for the backend and CustomTkinter for the interface.
    \item Built backend features and structured the API layer for communication between frontend and backend.
\end{itemize}}
\newpage
\section{Technical Skills}
\stackitem{Programming Languages}{Python, Java, JavaScript, C/C++, Dart}
\skilldivider
\stackitem{Backend and Web Technologies}{Node.js, Spring Boot, React, REST API, Socket.io}
\skilldivider
\stackitem{Machine Learning and Data Analysis}{Pandas, NumPy, LightGBM, XGBoost, CatBoost, TabNet, YOLO}
\skilldivider
\stackitem{Databases and Developer Tools}{MongoDB, Firebase, SQL, Git, Google Sheets API}
\skilldivider
\stackitem{Scripting and Command-Line Tools}{Bash/Shell Scripting, WSL}

\section{Languages}
\cvitemwithcomment{Turkish}{Native}{}
\cvitemwithcomment{English}{B2}{}
\cvitemwithcomment{German}{A2--B1}{}

\section{Certificates}
\cvitem{2026}{Artificial Intelligence Specialization Program - Foundation Training Certificate (Republic of T\"urkiye Ministry of Industry and Technology)}
\cvitem{2026}{Introduction to Open Source License Compliance Management (The Linux Foundation)}
\cvitem{2026}{Model Context Protocol: Advanced Topics (Anthropic)}
\cvitem{2026}{AI Fluency: Framework \& Foundations (Anthropic)}
\cvitem{2026}{Claude 101 (Anthropic)}
\cvitem{2025}{Google AI and Technology Academy Graduation Certificate}
\cvitem{2025}{Bootcamp Finalist Certificate of Achievement (Google AI \& Technology Academy)}
\cvitem{2025}{Google Advanced Data Analytics (Coursera)}
\cvitem{2025}{Foundations of Data Science (Google)}
\cvitem{2025}{The Nuts and Bolts of Machine Learning (Google)}
\cvitem{2025}{Regression Analysis: Simplify Complex Data Relationships (Google)}
\cvitem{2025}{Go Beyond the Numbers: Translate Data into Insights (Google)}
\cvitem{2025}{The Power of Statistics (Google)}
\cvitem{2025}{Get Started with Python (Google)}
\cvitem{2025}{Foundations of Project Management (Google)}
\cvitem{2022}{Deutsches Sprachdiplom 1 (Kultusministerkonferenz)}

\end{document}
